% 提案の手法（感情ごとの信頼度の表で声・顔・言葉の答えをまとめる）の図解と結果．2026-11-14 に組み直した（前の版は git の eee1f0a）．
% 作り方：.venv/bin/python DAI/docs/tools_pending/pdf_figs_proposal.py（結果の図表を結果ファイルから書き出す）→ cd DAI/docs && lualatex 説明_実験16の案_図解_2026-11-13.tex（2 回）
\documentclass[a4paper,11pt]{ltjsarticle}
\usepackage[margin=20mm]{geometry}
\usepackage{booktabs}
\usepackage{longtable}
\usepackage{array}
\usepackage{float}
\usepackage{xcolor}
\usepackage{amssymb}
\usepackage{tikz}
\usetikzlibrary{arrows.meta,positioning,fit,calc,patterns,shapes.geometric,backgrounds}
\usepackage{tcolorbox}
\usepackage{enumitem}
\usepackage[hidelinks]{hyperref}
\setlist{itemsep=2pt,topsep=4pt}

\definecolor{cA}{HTML}{2A6FB0}   % 声
\definecolor{cV}{HTML}{D2691E}   % 顔
\definecolor{cT}{HTML}{3A8F3A}   % 言葉
\definecolor{cSrc}{HTML}{4C78A8} % ソースのラベル付きデータ
\definecolor{cK}{HTML}{F28E2B}   % ターゲットの K 本
\definecolor{cU}{HTML}{BBBBBB}   % ターゲットのラベル無し
\definecolor{cTest}{HTML}{D62728}% テスト
\definecolor{cClf}{HTML}{59A14F} % 分類器
\definecolor{cRel}{HTML}{EDC948} % 信頼度の表
\definecolor{cMeta}{HTML}{9467BD}% メタ学習で決める数
\definecolor{cE16}{HTML}{1F5F8B}  % 実験 16 の印
\definecolor{cE17}{HTML}{8B3A62}  % 実験 17 の印
\definecolor{cE18}{HTML}{5A6B2F}  % 実験 18 の印
\definecolor{cBar1}{HTML}{4C78A8}
\definecolor{cBar2}{HTML}{F28E2B}
\definecolor{cBar3}{HTML}{9D9D9D}

\usepackage{graphicx}
\usepackage{caption}
\captionsetup{font=small,labelfont=bf}
\newtcolorbox{waku}[1]{colback=black!3,colframe=black!55,boxrule=0.6pt,arc=1mm,title={#1},fonttitle=\bfseries}
\newtcolorbox{chui}{colback=cTest!6,colframe=cTest!70,boxrule=0.6pt,arc=1mm}
\newcommand{\chip}[2]{\tikz[baseline=-0.6ex]\node[fill=#1,text=white,rounded corners=1.5pt,inner sep=2.2pt,font=\scriptsize\bfseries]{#2};}
\newcommand{\chipk}[2]{\tikz[baseline=-0.6ex]\node[fill=#1,text=black,rounded corners=1.5pt,inner sep=2.2pt,font=\scriptsize\bfseries]{#2};}
% 図を大きくする：1.22 倍にし，本文の幅を超えるなら幅いっぱいに収める
\usepackage{adjustbox}
\newcommand{\bigfigopts}{scale=1.22,max width=\linewidth}

\tikzset{
  box/.style={draw,rounded corners=2pt,align=center,inner sep=4pt,font=\small},
  clf/.style={draw=cClf!70!black,fill=cClf!18,rounded corners=2pt,align=center,inner sep=3pt,font=\scriptsize,minimum width=13mm,minimum height=7mm},
  arr/.style={-{Stealth[length=2.2mm]},thick},
  corp/.style={draw,rounded corners=3pt,minimum width=30mm,minimum height=12mm,align=center,font=\small},
}

\title{提案の手法の図解と結果\\[4pt]
\large 感情ごとの信頼度の表で，声・顔・言葉の分類器の答えをまとめる\\
\large （水準はターゲットの K 本で，相性はソースで測る）}
\author{MER-DA（DAI の作業場）}
\date{2026 年 11 月 14 日 更新}

\begin{document}
\maketitle

\begin{waku}{この資料の読み方}
\begin{center}
\small
\begin{tabular}{>{\raggedright}p{2.2cm}p{12.6cm}}
\toprule
節 & 中身 \\
\midrule
1〜3 & 何を解きたいか，使う部品，信頼度とは何か \\
4〜8 & 提案の手法の手順（4 つのステップ） \\
9 & 下調べの結果（4 コーパス） \\
10 & コーパスを 6 つに増やした結果 \\
11 & 分かったことと，まだ分からないこと \\
12 & いま計算している評価の回 \\
13・14 & 用語，結果ファイルの名前との対応 \\
\bottomrule
\end{tabular}
\end{center}
\begin{itemize}
  \item 「例」と書いた数は説明用に作った数である．9〜11 節の数は，結果ファイルから台本で書き出した．
  \item 9 節の結果は照合を通している．10 節の結果は出力の点検を通したが，結果の文の照合はこれからである．
  \item この資料は提案の手法だけを扱う．下調べで比べたほかの版と，脇に置いた案（近いソースの発話を引く案）は外した．全部の版の表は \texttt{DAI/docs/実験16・17\_下調べ\_結果の表\_K5・10・25\_2026-11-14.md} と \texttt{実験18\_下調べ\_結果の表\_2026-11-14.md} にある．
\end{itemize}
\end{waku}

\noindent 色の意味：
\chip{cSrc}{ソースのラベル付きデータ}\
\chip{cK}{ターゲットの K 本}\
\chipk{cU}{ターゲットのラベル無し}\
\chip{cTest}{テスト}\
\chip{cClf}{分類器}\
\chipk{cRel}{信頼度の表}\
\chip{cMeta}{別のコーパスで決める数}

%======================================================================
\section{何を解きたいか}

人が話している短い動画（発話）から，話し手の感情（Neutral，Happy，Sad，Angry の 4 つ）を当てる．
感情の正解（ラベル）が全部付いたデータの集まり（コーパス）が 3 つある．これを\textbf{ソース}と呼ぶ．
新しいコーパス（\textbf{ターゲット}）には，感情ごとに K 本（5 本，10 本か 25 本）しかラベルが無い．

\begin{figure}[H]
\centering
\begin{adjustbox}{scale=1.22,max width=\linewidth}\begin{tikzpicture}
  % ソース
  \foreach \nm/\y in {IEMOCAP/0,MELD/-1.5,MOSEI/-3.0}{
    \node[corp] (s\y) at (0,\y) {\nm};
    \fill[cSrc] (-1.5,\y-0.85) rectangle (1.5,\y-0.65);
    \node[font=\scriptsize,anchor=west] at (1.6,\y-0.75) {ラベル：全部ある};
  }
  \node[font=\small\bfseries] at (0,0.95) {ソース（3 コーパス）};
  % ターゲット
  \node[corp,minimum width=44mm] (t) at (8.2,-1.5) {AFEW（新しいコーパス）};
  \fill[cK] (6.0,-2.35) rectangle (6.5,-2.15);
  \fill[cU] (6.5,-2.35) rectangle (9.2,-2.15);
  \draw[cTest,very thick] (9.3,-2.35) rectangle (10.4,-2.15);
  \node[font=\scriptsize,cK!70!black,anchor=north] at (6.25,-2.4) {K 本};
  \node[font=\scriptsize,anchor=north] at (7.85,-2.4) {ラベル無し};
  \node[font=\scriptsize,cTest,anchor=north] at (9.85,-2.4) {テスト};
  \node[font=\small\bfseries] at (8.2,-0.55) {ターゲット};
  \draw[arr] (2.9,-1.5) -- node[above,font=\scriptsize]{ソースで学んだことを} node[below,font=\scriptsize]{持ち込む} (5.9,-1.5);
  \node[font=\scriptsize,align=left,anchor=west] at (5.9,-3.6) {テストの発話の正解は，成績を測るときにだけ使う．\\テストの話者は，K 本にもラベル無しの発話にも入れない．};
\end{tikzpicture}\end{adjustbox}
\caption{問題の設定．ソースはラベルが豊富，ターゲットはラベルが少ない．図は AFEW をターゲットにした場合で，4 つのコーパスを順にターゲットにする．}
\end{figure}

成績は\textbf{UAR}で測る．UAR は，4 つの感情それぞれの正解率を平均した値である（でたらめなら 25 点）．

%======================================================================
\section{使う部品}

発話からは声・顔・言葉の 3 種類の情報（\textbf{モダリティ}）が取れる．それぞれを学習済みのネットワーク（エンコーダ．中身は変えない）で数の並びに変え，その上に線形の\textbf{分類器}を置く．

\begin{figure}[H]
\centering
\begin{adjustbox}{scale=1.22,max width=\linewidth}\begin{tikzpicture}
  \node[font=\small\bfseries] at (2.6,1.0) {ソースごとの分類器（3 × 3 = 9 個）};
  \foreach \m/\x/\c in {声/1.2/cA,顔/2.6/cV,言葉/4.0/cT}{ \node[font=\small,\c] at (\x,0.45) {\m}; }
  \foreach \s/\y in {IEMOCAP/-0.2,MELD/-1.1,MOSEI/-2.0}{
    \node[font=\scriptsize,anchor=east] at (0.35,\y) {\s};
    \foreach \x in {1.2,2.6,4.0}{ \node[clf] at (\x,\y) {分類器}; }
  }
  \node[font=\scriptsize,align=center] at (2.6,-2.85) {ソースの全ラベルで学ぶ};
  % 微調整
  \node[font=\small\bfseries] at (9.3,1.0) {ターゲットで微調整した分類器（3 個）};
  \foreach \m/\x/\c in {声/7.9/cA,顔/9.3/cV,言葉/10.7/cT}{
    \node[font=\small,\c] at (\x,0.45) {\m};
    \node[clf,fill=cK!25,draw=cK!70!black] (f\x) at (\x,-0.2) {分類器};
  }
  \node[font=\scriptsize,align=center] at (9.3,-1.05) {ソースの分類器（モダリティごとに 3 個）の\\平均から始めて，ターゲットの K 本で学び直す};
  \draw[arr] (4.9,-1.1) -- (7.0,-0.3);
\end{tikzpicture}\end{adjustbox}
\caption{部品．分類器は，発話ごとに 4 つの感情の確率を出す．}
\end{figure}

\noindent 解きたいのは，\textbf{これらの分類器の答えを，発話ごとに，どれをどの程度信頼してまとめるか}である．今までの方法は，微調整した 3 つの分類器の確率を同じ強さで掛け合わせる．これを\textbf{積}と呼ぶ．

\subsection*{分類器のつながり}

\begin{figure}[H]
\centering
\begin{adjustbox}{scale=1.22,max width=\linewidth}\begin{tikzpicture}[
  st/.style={draw,rounded corners=2pt,align=center,font=\scriptsize,inner sep=3pt,minimum height=12mm},
]
  \node[st,fill=cSrc!18,minimum width=30mm] (a) at (0,0) {\textbf{(1) ソースの分類器}\\コーパス × モダリティごと\\（ソースごと × 声・顔・言葉，\\と融合）};
  \node[st,fill=cClf!15,minimum width=26mm] (b) at (3.7,0) {\textbf{(2) 出発点}\\モダリティごとに，\\ソースの分類器を\\平均する};
  \node[st,fill=cK!18,minimum width=28mm] (c) at (7.2,0) {\textbf{(3) K 本だけの融合}\\融合の出発点から\\K 本で学ぶ};
  \node[st,fill=cK!28,minimum width=30mm] (d) at (10.9,0) {\textbf{(4) 微調整した}\\\textbf{声・顔・言葉}\\出発点から，K 本と\\(3) の疑似ラベルで学ぶ};
  \draw[arr] (a) -- (b); \draw[arr] (b) -- (c); \draw[arr] (c) -- (d);
  \draw[arr] (b.south) |- (9.0,-1.45) -| (d.south);
  \node[st,fill=black!5,minimum width=36mm] (e1) at (7.6,-3.0) {\textbf{積（今までの方法）}\\3 つの確率を同じ重みで掛ける};
  \node[st,fill=cRel!40,minimum width=40mm] (e2) at (12.2,-3.0) {\textbf{提案の手法}\\信頼度の表から作った重みで\\3 つの点を足す（8 節）};
  \draw[arr] (d.south) ++(-0.4,0) -- (e1.north);
  \draw[arr] (d.south) ++(0.4,0) -- (e2.north);
  \node[st,fill=cTest!8,minimum width=26mm] (f) at (9.9,-4.75) {答え（感情）};
  \draw[arr] (e1) -- (f); \draw[arr] (e2) -- (f);
  \node[font=\scriptsize,align=left,anchor=west] at (-1.5,-3.0) {提案の手法が変えるのは，\\「3 つの答えをどうまとめるか」だけである．\\まとめた答えは，融合の分類器を学び直すときの\\疑似ラベルにも使える（10 節）．};
\end{tikzpicture}\end{adjustbox}
\caption{分類器のつながり．ソースの分類器はコーパス × モダリティごとにあるが，最後にまとめるのは (4) の 3 つ（声・顔・言葉）である．ソースの分類器は (2) の出発点として (4) に入る．}
\end{figure}


%======================================================================
\section{信頼度とは何か}

\textbf{信頼度}は，分類器ごと・感情ごとに 1 つずつ持つ数で，「本当がその感情のとき，その分類器をどれだけ頼ってよいか」を表す．

\begin{figure}[H]
\centering
\begin{adjustbox}{scale=1.22,max width=\linewidth}\begin{tikzpicture}
  % 表
  \node[font=\small\bfseries] at (2.0,1.25) {信頼度の表（例）};
  \foreach \e/\x in {Neutral/1.1,Happy/2.3,Sad/3.5,Angry/4.7}{ \node[font=\scriptsize] at (\x,0.65) {\e}; }
  \node[font=\small,cV,anchor=east] at (0.45,0) {顔};
  \node[font=\small,cA,anchor=east] at (0.45,-0.75) {声};
  \node[font=\small,cT,anchor=east] at (0.45,-1.5) {言葉};
  \foreach \v/\x/\y/\d in {0.5/1.1/0/40,0.9/2.3/0/100,0.6/3.5/0/55,0.8/4.7/0/85,
                          0.5/1.1/-0.75/40,0.4/2.3/-0.75/25,0.8/3.5/-0.75/85,0.7/4.7/-0.75/70,
                          0.6/1.1/-1.5/55,0.5/2.3/-1.5/40,0.5/3.5/-1.5/40,0.6/4.7/-1.5/55}{
    \node[fill=cRel!\d,minimum width=11.4mm,minimum height=6.9mm,inner sep=0pt,font=\small] at (\x,\y) {\v};
  }
  \node[font=\scriptsize,align=left,anchor=west] at (0.2,-2.25) {色が濃いほど頼れる．\\顔は Happy に強く，声は Sad に強い．};
  % 分解
  \begin{scope}[xshift=7.6cm,yshift=-1.75cm]
    \node[font=\small\bfseries] at (2.6,3.45) {顔の行を 2 つに分ける};
    \draw[->] (0,0) -- (0,2.95);
    \node[font=\scriptsize,anchor=south west] at (-0.05,2.95) {信頼度};
    \draw (0,0) -- (5.2,0);
    \foreach \v/\x/\e in {0.5/0.75/Neutral,0.9/2.05/Happy,0.6/3.35/Sad,0.8/4.65/Angry}{
      \fill[cRel!80] (\x-0.38,0) rectangle (\x+0.38,\v*3.0);
      \node[font=\scriptsize] at (\x,-0.24) {\e};
    }
    \draw[cTest,thick,dashed] (0,0.7*3.0) -- (5.2,0.7*3.0);
    \node[font=\scriptsize,cTest,anchor=west,align=left] at (5.25,0.7*3.0) {\textbf{水準} 0.7\\（行の平均）};
    \draw[<->,thick] (2.05,0.7*3.0) -- (2.05,0.9*3.0);
    \node[font=\scriptsize,anchor=west] at (2.5,0.82*3.0) {\textbf{相性} $+0.2$};
    \draw[<->,thick] (0.75,0.5*3.0) -- (0.75,0.7*3.0);
    \node[font=\scriptsize,anchor=west] at (1.17,0.58*3.0) {$-0.2$};
  \end{scope}
\end{tikzpicture}\end{adjustbox}
\caption{左：信頼度の表の例．右：1 つの行を「水準」（その分類器が全体にどれだけ強いか）と「相性」（感情ごとの得意と不得意）に分ける．}
\end{figure}

\begin{itemize}
  \item \textbf{水準}はコーパスで変わる（記録：AFEW は顔，IEMOCAP は声，MELD は言葉が最も強い）．ソースからは分からないので，ターゲットの K 本で測る．
  \item \textbf{相性}は，コーパスをまたいで似ている可能性がある（記録：顔の分類器が最も当てた感情は 4 コーパスとも Happy）．数が多く K 本では測れないので，ソースで測る．相性がどこまで持ち越せるかは，7.3 節と 11 節で数を示す．
\end{itemize}

%======================================================================
\section{全体の流れ}

\begin{figure}[H]
\centering
\begin{adjustbox}{scale=1.22,max width=\linewidth}\begin{tikzpicture}[
  st/.style={draw,rounded corners=3pt,minimum width=33mm,minimum height=31mm,align=center,font=\small,inner sep=2pt},
]
  \node[st,fill=cSrc!10] (a) at (0,0) {\textbf{① 学習時}\\[2pt]ソースで分類器を作る\\[2pt]ソースで信頼度を測る\\{\scriptsize（中で測る → 外で測り直す）}};
  \node[st,fill=cMeta!10] (b) at (3.9,0) {\textbf{② メタ学習}\\[2pt]ソースの 1 つを\\仮のターゲットにして，\\「決め方の数」を決める};
  \node[st,fill=cK!12] (c) at (7.8,0) {\textbf{③ ターゲットで適応}\\[2pt]K 本で微調整する\\[2pt]K 本で水準を測り，\\表を組み立てる};
  \node[st,fill=cTest!8] (d) at (11.7,0) {\textbf{④ テスト}\\[2pt]発話ごとに，\\信頼度の表を使って\\分類器の答えをまとめる};
  \draw[arr] (a) -- (b); \draw[arr] (b) -- (c); \draw[arr] (c) -- (d);
  % 使うデータ
  \node[font=\scriptsize,anchor=north] at (0,-1.7) {\chip{cSrc}{ソース}};
  \node[font=\scriptsize,anchor=north] at (3.9,-1.7) {\chip{cSrc}{ソース}};
  \node[font=\scriptsize,anchor=north] at (7.8,-1.7) {\chip{cK}{K 本}\ \chipk{cU}{ラベル無し}};
  \node[font=\scriptsize,anchor=north] at (11.7,-1.7) {\chip{cTest}{テストの発話}};
  \node[font=\scriptsize,anchor=east] at (-1.75,-1.92) {使うデータ};
  % できるもの
  \node[font=\scriptsize,anchor=north,align=center] at (0,-2.4) {\chip{cClf}{分類器 9 個}\\[2pt]\chipk{cRel}{相性の表}};
  \node[font=\scriptsize,anchor=north,align=center] at (3.9,-2.4) {\chip{cMeta}{組み立て方の数 3 つ}};
  \node[font=\scriptsize,anchor=north,align=center] at (7.8,-2.4) {\chip{cClf}{微調整した分類器}\\[2pt]\chipk{cRel}{信頼度の表（完成）}};
  \node[font=\scriptsize,anchor=north,align=center] at (11.7,-2.4) {感情の答え};
  \node[font=\scriptsize,anchor=east] at (-1.75,-2.62) {できるもの};
\end{tikzpicture}\end{adjustbox}
\caption{全体の流れ．①と②はソースだけで済ませる．ターゲットのラベルを使うのは③の K 本だけである．}
\end{figure}

\begin{chui}
\textbf{テストの正解はどこにも入らない．}①〜④のどのステップでも，テストの発話の正解は使わない．正解は，④で出した答えの成績（UAR）を測るときにだけ使う．
\end{chui}

%======================================================================
\section{① 学習時：ソースで信頼度の表を作る}

この節では，ターゲットのラベルを 1 本も使わずに，ソースだけで「信頼度の表」の材料を作る手順を説明する．
信頼度の表は，分類器 3 つ（声・顔・言葉）× 感情 4 つの 12 個の数である（3 節）．

\subsection*{5.1 1 回の計算に出てくる 4 つのコーパスの役}

下調べでは，4 つのコーパスに次の 3 つの役を割り当てて 1 回の計算をする．

\begin{figure}[H]
\centering
\begin{adjustbox}{scale=1.22,max width=\linewidth}\begin{tikzpicture}
  \node[draw=black!30,fill=black!6,rounded corners=3pt,minimum width=30mm,minimum height=17mm,align=center,font=\small,text=black!55] at (0,0) {MOSEI\\[2pt]{\scriptsize\textbf{評価のターゲット}}\\{\scriptsize 触らない（本番用に残す）}};
  \node[draw=cK!80!black,line width=1.4pt,fill=cK!20,rounded corners=3pt,minimum width=36mm,minimum height=17mm,align=center,font=\small] at (4.3,0) {AFEW\\[2pt]{\scriptsize\textbf{仮のターゲット}}\\{\scriptsize K 本・ラベル無し・テスト}};
  \node[draw=cSrc!80!black,fill=cSrc!20,rounded corners=3pt,minimum width=30mm,minimum height=17mm,align=center,font=\small] at (8.6,0) {IEMOCAP\\[2pt]{\scriptsize\textbf{ソース 1}}\\{\scriptsize ラベルを全部使う}};
  \node[draw=cSrc!80!black,fill=cSrc!20,rounded corners=3pt,minimum width=30mm,minimum height=17mm,align=center,font=\small] at (12.2,0) {MELD\\[2pt]{\scriptsize\textbf{ソース 2}}\\{\scriptsize ラベルを全部使う}};
  % 仮のターゲットの中身
  \fill[cK] (2.6,-1.35) rectangle (3.2,-1.15);
  \fill[cU] (3.2,-1.35) rectangle (5.1,-1.15);
  \draw[cTest,very thick] (5.2,-1.35) rectangle (6.0,-1.15);
  \node[font=\tiny,anchor=north] at (2.9,-1.4) {K 本};
  \node[font=\tiny,anchor=north] at (4.15,-1.4) {ラベル無し};
  \node[font=\tiny,anchor=north,cTest] at (5.6,-1.4) {テスト};
\end{tikzpicture}\end{adjustbox}
\caption{1 回の計算の役の例．評価のターゲットは，本番で一度だけ試すために残し，この下調べでは中身を一切見ない．仮のターゲットは全ラベルを持つので，「K 本しかない」状況を作って答え合わせができる．}
\end{figure}

\begin{itemize}
  \item 評価のターゲットの選び方が 4 通り，残り 3 つから仮のターゲットの選び方が 3 通りで，組み合わせは 12 通りある．ソースは残りの 2 つである．
  \item 仮のターゲットの中の話者の分け方（5 通り）と K 本の引き方（2 通り）も替えて，計 12 × 5 × 2 = 120 回計算する．
  \item たとえば AFEW を仮のターゲットにする計算は，ソースの組（IEMOCAP と MELD，IEMOCAP と MOSEI，MELD と MOSEI）を替えて 3 通りある．AFEW の結果はこの 3 通りの平均である．
\end{itemize}

\subsection*{5.2 ソースの分類器ごとに，中で測り，外で測り直す}

ソースの分類器は，モダリティごとにソースの数だけある（この例では声の分類器が IEMOCAP の分と MELD の分の 2 つ）．
1 つの分類器について，信頼度の行（感情 4 つ）を 2 通りに測る．

\begin{figure}[H]
\centering
\begin{adjustbox}{scale=1.22,max width=\linewidth}\begin{tikzpicture}
  % IEMOCAP の声の分類器
  \node[clf,minimum width=26mm] (ci) at (0,1.2) {IEMOCAP の\\声の分類器};
  \node[clf,minimum width=26mm] (cm) at (0,-1.6) {MELD の\\声の分類器};
  % 中
  \node[draw=cSrc!80!black,fill=cSrc!12,rounded corners=2pt,font=\scriptsize,align=center,minimum width=34mm] (hi) at (4.6,1.9) {\textbf{中で測る}：IEMOCAP の中の\\話者で分けた交差検証};
  \node[draw=cSrc!80!black,fill=cSrc!12,rounded corners=2pt,font=\scriptsize,align=center,minimum width=34mm] (ai) at (4.6,0.5) {\textbf{外で測る}：MELD の\\全発話に当てる};
  \node[draw=cSrc!80!black,fill=cSrc!12,rounded corners=2pt,font=\scriptsize,align=center,minimum width=34mm] (hm) at (4.6,-0.9) {\textbf{中で測る}：MELD の中の\\話者で分けた交差検証};
  \node[draw=cSrc!80!black,fill=cSrc!12,rounded corners=2pt,font=\scriptsize,align=center,minimum width=34mm] (am) at (4.6,-2.3) {\textbf{外で測る}：IEMOCAP の\\全発話に当てる};
  \draw[arr] (ci) -- (hi.west); \draw[arr] (ci) -- (ai.west); \draw[arr] (cm) -- (hm.west); \draw[arr] (cm) -- (am.west);
  % 行の値（例）
  \foreach \e/\x in {N/7.75,H/8.55,S/9.35,A/10.15}{ \node[font=\scriptsize] at (\x,2.65) {\e}; }
  \node[font=\scriptsize,anchor=west] at (10.6,2.65) {合計};
  \foreach \vals/\y/\s in {{0.8,0.4,1.2,0.8}/1.9/3.2,{0.6,0.3,0.7,0.6}/0.5/2.2,{0.8,0.4,0.8,0.8}/-0.9/2.8,{0.6,0.4,0.7,0.6}/-2.3/2.3}{
    \foreach \v [count=\j from 0] in \vals { \node[draw,fill=cRel!50,minimum width=7.5mm,minimum height=5.5mm,inner sep=0pt,font=\scriptsize] at (7.75+0.8*\j,\y) {\v}; }
    \node[font=\scriptsize,anchor=west] at (10.6,\y) {\s};
  }
\end{tikzpicture}\end{adjustbox}
\caption{ソースの声の分類器 2 つの，中と外の信頼度（例）．N・H・S・A は Neutral・Happy・Sad・Angry．中で測った値は「同じコーパスなら，これくらい頼れる」，外で測った値は「別のコーパスに持ち出すと，これくらいに落ちる」を表す．}
\end{figure}

\noindent\textbf{残り方}：外に持ち出すとどれだけ残るかを，モダリティごとに 1 つの数にまとめる．
\begin{center}
\small
残り方（声）＝ 外で測った値の合計 ÷ 中で測った値の合計 ＝（2.2 ＋ 2.3）÷（3.2 ＋ 2.8）＝ 4.5 ÷ 6.0 ＝ \textbf{0.75}
\end{center}
外での値は相手のコーパスでしか測れず，ターゲットではまだ測れない．そこで「ターゲットに持ち出しても，別のソースに持ち出したときと同じくらい残る」と置き，\textbf{中で測った値 × 残り方}を，ソースで見込むターゲットでの信頼度にする．

\subsection*{5.3 微調整した分類器の行にする}

ターゲットで使う声の分類器は，ソースの声の分類器 2 つの平均から始めて，K 本で微調整したものである（2 節）．そこで，その行も 2 つの平均にする．

\begin{figure}[H]
\centering
\begin{adjustbox}{scale=1.22,max width=\linewidth}\begin{tikzpicture}
  \foreach \e/\x in {N/1.0,H/1.8,S/2.6,A/3.4}{ \node[font=\scriptsize] at (\x,1.55) {\e}; }
  \node[font=\scriptsize,anchor=east] at (0.5,0.9) {IEMOCAP：中 × 0.75};
  \node[font=\scriptsize,anchor=east] at (0.5,0.2) {MELD：中 × 0.75};
  \foreach \vals/\y in {{0.60,0.30,0.90,0.60}/0.9,{0.60,0.30,0.60,0.60}/0.2}{
    \foreach \v [count=\j from 0] in \vals { \node[draw,fill=cSrc!30,minimum width=7.5mm,minimum height=5.5mm,inner sep=0pt,font=\scriptsize] at (1.0+0.8*\j,\y) {\v}; }
  }
  \draw[arr] (3.95,0.55) -- node[above,font=\scriptsize]{平均} (5.2,0.55);
  \foreach \v [count=\j from 0] in {0.60,0.30,0.75,0.60} { \node[draw,fill=cSrc!45,minimum width=7.5mm,minimum height=5.5mm,inner sep=0pt,font=\scriptsize] at (5.7+0.8*\j,0.55) {\v}; }
  \foreach \e/\x in {N/5.7,H/6.5,S/7.3,A/8.1}{ \node[font=\scriptsize] at (\x,1.2) {\e}; }
  \node[font=\scriptsize,anchor=west] at (8.6,0.55) {← 声の行（ソースで見込んだ値）};
  % 水準と相性に分ける
  \draw[arr] (6.9,0.1) -- (6.9,-0.55);
  \node[draw,fill=cRel!35,rounded corners=2pt,font=\scriptsize,align=left,anchor=north] at (6.9,-0.6) {水準（行の平均）0.56\\相性（平均との差）：N +0.04，H −0.26，S +0.19，A +0.04};
\end{tikzpicture}\end{adjustbox}
\caption{微調整した声の分類器の，ソースで見込んだ行（例）．顔と言葉の行も同じ手順で作る．こうして 3 × 4 の表ができる．}
\end{figure}

\begin{chui}
\textbf{ここまでにターゲットは出てこない．}5 節の表は，ソース 2 つのラベルだけで作る．仮のターゲットの K 本・ラベル無しの発話・テストは，どれも使わない．この表をそのまま使う版を「ソースだけで測った表」と呼び，9 節で提案と並べる．
\end{chui}

%======================================================================
\section{② メタ学習：ターゲットでの組み立て方を，別のコーパスで決める}

\subsection*{6.1 何を決めるか}

ターゲットで信頼度の表を組み立てるとき（7 節），前もって決めておく数が 3 つある．

\begin{center}
\small
\begin{tabular}{>{\raggedright}p{4.3cm}>{\raggedright}p{7.4cm}p{3.3cm}}
\toprule
決める数 & 意味 & 候補 \\
\midrule
戻す強さ n0 & K 本で測った水準を，どれだけ信じるか（7.2 節で図にする） & 0，20，80，∞（ラベルの本数） \\
戻し先 & 水準を，どこへ向けて戻すか & 3 つの平均（等しい重み），ソースで見込んだ水準 \\
相性を使うか β & 5 節の相性を表に入れるか & 入れない（0），入れる（1） \\
\bottomrule
\end{tabular}
\end{center}
組み合わせは 4 × 2 × 2 ＝ 16 通りである（n0 ＝ 0 では戻し先が効かないので，中身が違うのは 14 通り）．

\subsection*{6.2 なぜ別のコーパスで決めるか}

16 通りのうちどれが良いかは，テストで成績を測れば分かる．しかし，同じコーパスのテストで選ぶと，そのテストに合わせた選び方になり，成績が甘く出る．そこで，\textbf{成績を測るコーパスを一切含まない計算で選ぶ}．

\begin{figure}[H]
\centering
\begin{adjustbox}{scale=1.22,max width=\linewidth}\begin{tikzpicture}[
  cc/.style={draw,rounded corners=2pt,minimum width=17mm,minimum height=6.5mm,font=\scriptsize,inner sep=1pt},
]
  \node[font=\small\bfseries,anchor=west] at (-1.8,2.0) {AFEW の成績に使う設定を選ぶとき};
  \node[font=\scriptsize,anchor=west] at (-1.8,1.45) {AFEW を「評価のターゲット」（触らない側）に置いた 30 回だけを使う};
  \node[font=\scriptsize\bfseries] at (0.35,0.85) {仮のターゲット};
  \node[font=\scriptsize\bfseries] at (3.6,0.85) {ソース};
  \node[font=\scriptsize\bfseries] at (6.4,0.85) {触らない};
  \foreach \t/\sa/\sb/\y in {IEMOCAP/MOSEI/MELD/0.2,MOSEI/IEMOCAP/MELD/-0.6,MELD/IEMOCAP/MOSEI/-1.4}{
    \node[cc,draw=cK!80!black,fill=cK!25] at (0.35,\y) {\t};
    \node[cc,draw=cSrc!80!black,fill=cSrc!20] at (2.75,\y) {\sa};
    \node[cc,draw=cSrc!80!black,fill=cSrc!20] at (4.45,\y) {\sb};
    \node[cc,draw=black!30,fill=black!6,text=black!50] at (6.4,\y) {AFEW};
  }
  \node[font=\scriptsize,anchor=west] at (-1.8,-2.1) {それぞれ 分割 5 × 種 2 ＝ 10 回．計 30 回};
  % 手順
  \begin{scope}[xshift=10.9cm,yshift=1.3cm]
    \node[box,fill=cK!15,minimum width=52mm] (s1) at (0,0) {30 回のそれぞれで，16 通りの設定で\\表を組み立てて答えをまとめる};
    \node[box,fill=cTest!8,minimum width=56mm] (s2) at (0,-1.2) {仮のターゲットのテストで UAR を測り，\\設定ごとに 30 回の平均を出す};
    \node[box,fill=cMeta!15,minimum width=56mm] (s3) at (0,-2.4) {平均が最も高い設定を 1 つ選ぶ};
    \node[box,fill=cK!25,minimum width=56mm] (s4) at (0,-3.6) {その設定を，AFEW を仮のターゲットに\\した計算（AFEW の成績）に使う};
    \draw[arr] (s1) -- (s2); \draw[arr] (s2) -- (s3); \draw[arr] (s3) -- (s4);
  \end{scope}
\end{tikzpicture}\end{adjustbox}
\caption{メタ学習の中身．「AFEW のための設定」を，AFEW が仮のターゲットにもソースにも出てこない 30 回で選ぶ．IEMOCAP・MOSEI・MELD の成績に使う設定も，それぞれ同じやり方で選ぶ．選ぶのは K ごとである．}
\end{figure}

\begin{itemize}
  \item 「メタ学習」と呼んでいるが，実際にしているのは「16 通りの中から 1 つを，別のコーパスでの成績で選ぶ」ことである．
  \item こうすると，AFEW のテストの正解は，AFEW の設定の選び方に入らない．
  \item 選ぶ計算の中にも，ソース 2 つと仮のターゲット 1 つがあり，評価する計算と同じ形をしている．
\end{itemize}

%======================================================================
\section{③ ターゲットで適応：K 本で水準を測り，表を組み立てる}

ここで初めてターゲットの K 本を使う．K = 5 なら 4 感情 × 5 本 ＝ 20 本である．

\subsection*{7.1 K 本で水準を測る}

\begin{figure}[H]
\centering
\begin{adjustbox}{scale=1.22,max width=\linewidth}\begin{tikzpicture}
  \node[font=\small\bfseries,anchor=west] at (-0.3,1.6) {20 本を 5 組に分ける（1 組 ＝ 各感情 1 本ずつ，計 4 本）};
  \foreach \g in {0,1,2,3,4}{
    \foreach \e/\c in {0/cV,1/cA,2/cT,3/cK}{ \fill[cK!70] (\g*1.6+\e*0.32,0.8) rectangle (\g*1.6+\e*0.32+0.27,1.05); }
    \node[font=\tiny] at (\g*1.6+0.6,0.6) {組 \g};
  }
  \draw[cTest,thick] (-0.08,0.72) rectangle (1.3,1.13);
  \node[font=\scriptsize,align=left,anchor=west] at (8.2,0.92) {組 0 を抜いた 16 本で分類器を微調整し直し，\\抜いた組 0 の 4 本に点を付ける．組を替えて 5 回．};
  \draw[arr] (3.9,0.45) -- (3.9,-0.15);
  \node[box,fill=cRel!35,minimum width=110mm,align=left] at (5.3,-0.75) {20 本すべてに「その発話を使わずに学んだ分類器」の点が付く．\\これで分類器ごと・感情ごとの信頼度を測り，行の平均を K 本で測った\textbf{水準}にする．};
\end{tikzpicture}\end{adjustbox}
\caption{K 本で水準を測る手順．分類器は K 本を学習に使っているので，同じ発話でそのまま測ると強く見えすぎる．1 組ずつ抜いて測ることで，これを避ける．}
\end{figure}

20 本しかないので，測った水準はぶれる．このぶれを抑えるのが，次の「戻す」操作である．

\subsection*{7.2 水準を戻す（戻す強さ n0）}

\textbf{考え方}：K 本で測った水準に，「3 つの分類器はどれも同じくらい強い」と言う仮のラベルを n0 本ぶん混ぜて平均する．n0 が大きいほど，K 本の測定より「どれも同じ」を信じる．
\begin{center}
\small
使う水準 ＝（K 本の本数 × K 本で測った水準 ＋ n0 × 戻し先の水準）÷（K 本の本数 ＋ n0）
\end{center}
例：K 本で測った水準が 顔 0.70，声 0.48，言葉 0.42 だったとする．3 つの平均は 0.533 である．戻し先を「3 つの平均」，n0 ＝ 20 にすると，顔は（20 × 0.70 ＋ 20 × 0.533）÷ 40 ＝ 0.617 になる．

\begin{figure}[H]
\centering
\begin{adjustbox}{scale=1.22,max width=\linewidth}\begin{tikzpicture}
  \foreach \n/\ox/\fv/\av/\tv in {0/0/0.70/0.48/0.42,20/3.9/0.617/0.507/0.477,80/7.8/0.567/0.523/0.511,{$\infty$}/11.7/0.533/0.533/0.533}{
    \begin{scope}[xshift=\ox cm]
      \draw (0,0) -- (3.0,0);
      \draw[dashed,black!50] (0,0.533*3.2) -- (3.0,0.533*3.2);
      \foreach \v/\x/\c/\m in {\fv/0.5/cV/顔,\av/1.5/cA/声,\tv/2.5/cT/言葉}{
        \fill[\c] (\x-0.33,0) rectangle (\x+0.33,\v*3.2);
        \node[font=\tiny,fill=white,inner sep=0.5pt] at (\x,\v*3.2+0.17) {\pgfmathprintnumber[fixed,fixed zerofill,precision=2]{\v}};
        \node[font=\scriptsize,\c] at (\x,-0.22) {\m};
      }
      \node[font=\small\bfseries] at (1.5,2.75) {n0 ＝ \n};
    \end{scope}
  }
  \node[font=\scriptsize,anchor=north,align=center] at (1.5,-0.5) {K 本の値をそのまま使う};
  \node[font=\scriptsize,anchor=north,align=center] at (5.4,-0.5) {平均との差が半分になる};
  \node[font=\scriptsize,anchor=north,align=center] at (9.3,-0.5) {差が 5 分の 1 になる};
  \node[font=\scriptsize,anchor=north,align=center] at (13.2,-0.5) {3 つとも同じ\\（積と同じ重み）};
\end{tikzpicture}\end{adjustbox}
\caption{戻す強さ n0 による水準の変わり方（K = 5，ラベル 20 本の例）．点線は 3 つの平均 0.533．n0 はラベルの本数で数えるので，K が増えると同じ n0 でも K 本の測定が強く効く（K = 25 の 100 本なら，n0 ＝ 20 でも差は 6 分の 5 残る）．}
\end{figure}

\begin{itemize}
  \item \textbf{戻し先}は 2 通りある．「3 つの平均」は上の図のとおりである．「ソースで見込んだ水準」は，5 節の表の行の平均へ向けて戻す．
  \item n0 ＝ ∞ で相性を入れないと，表は全部同じ数になり，積と同じ答えになる．
\end{itemize}

\subsection*{7.3 水準に相性を足して表にする}

表の 1 行（1 つの分類器）は，次の 2 つに分けられる（3 節）．
\begin{itemize}
  \item \textbf{水準}：行の平均．1 つの数で，「この分類器は全体としてどれだけ当たるか」を表す．
  \item \textbf{相性}：行の各値から平均を引いた 4 つの数（足すと 0）．「この分類器は，どの感情が比較的得意で，どれが苦手か」を表す．
\end{itemize}
水準と相性は，別々のデータから取る．

\noindent\textbf{水準をソースで決められない理由}：「声・顔・言葉のどれがいちばん当たるか」は，コーパスごとに違う．下の表は，K = 5 で微調整した分類器の UAR（下調べの記録）である．
\begin{center}
\small
\begin{tabular}{lccc}
\toprule
コーパス & 声 & 顔 & 言葉 \\
\midrule
AFEW & 52.2 & \textbf{70.8} & 45.9 \\
IEMOCAP & \textbf{67.6} & 51.4 & 56.4 \\
MOSEI & \textbf{37.5} & 37.1 & 34.1 \\
MELD & 43.4 & 37.2 & \textbf{49.0} \\
\bottomrule
\end{tabular}
\end{center}
たとえば AFEW をターゲットにして，ソースが IEMOCAP と MELD のとき，ソースで測ると「声と言葉が強く，顔は弱い」と出る．しかし AFEW では顔がいちばん強い．このように，ソースで測った強さの順は，ターゲットでの順を当てない．
実際に，ソースで決めた表と K 本で決めた表で「いちばん強いモダリティ」を選び，全ラベルで測った正解と合った割合を比べると，次のようになった（K = 5）．
\begin{center}
\small
\begin{tabular}{lcccc}
\toprule
 & AFEW & IEMOCAP & MOSEI & MELD \\
\midrule
ソースで決めた表（5 節） & 0.33 & 0.00 & 0.37 & 0.00 \\
ターゲットの K 本（20 本）で決めた表 & 0.97 & 0.67 & 0.43 & 0.77 \\
\bottomrule
\end{tabular}
\end{center}

\noindent\textbf{水準なら 20 本で測れる理由}：水準は分類器 1 つにつき 1 つの数で，20 本を全部使って測る．相性は感情ごとの数なので，1 つの数に使えるのはその感情の 5 本だけである．5 本では，たまたま当たった・外れたで値が大きく動く．

\noindent\textbf{相性をソースで取る理由}：「同じ分類器の中で，どの感情が比較的得意か」は，コーパスをまたいでおおむね似ている（記録：顔の分類器が比較的よく当てた感情は 4 コーパスとも Happy．ソースで測った相性の形と，全ラベルで測った形の相関は 0.49〜0.77）．そこで相性は，ラベルの多いソースで測った値を使う．

\begin{figure}[H]
\centering
\begin{adjustbox}{scale=1.22,max width=\linewidth}\begin{tikzpicture}[c/.style={draw,minimum width=7.6mm,minimum height=5.4mm,inner sep=0pt,font=\scriptsize}]
  % (A) ソースで見込んだ表
  \node[font=\scriptsize\bfseries,align=center] at (1.5,2.35) {(A) ソースで見込んだ表\\（5 節）};
  \foreach \e/\x in {N/0.4,H/1.15,S/1.9,A/2.65}{ \node[font=\tiny] at (\x,1.75) {\e}; }
  \foreach \m/\c/\y/\vals in {顔/cV/1.3/{0.50,0.90,0.40,0.60},声/cA/0.7/{0.60,0.30,0.75,0.60},言葉/cT/0.1/{0.55,0.45,0.45,0.55}}{
    \node[font=\scriptsize,\c,anchor=east] at (0.0,\y) {\m};
    \foreach \v [count=\j from 0] in \vals { \node[c,fill=cSrc!35] at (0.4+0.75*\j,\y) {\v}; }
  }
  \draw[arr] (3.15,0.7) -- node[above,font=\tiny,align=center]{行ごとに\\分ける} (4.15,0.7);
  % (B) 水準（ソース）と相性
  \node[font=\scriptsize\bfseries,align=center] at (4.75,2.35) {水準\\（ソース）};
  \foreach \v/\y in {0.60/1.3,0.56/0.7,0.50/0.1}{ \node[c,fill=cSrc!20,text=black!45] at (4.75,\y) {\v}; }
  \node[font=\scriptsize] at (4.75,-0.5) {捨てる};
  \node[font=\scriptsize\bfseries,align=center] at (7.3,2.35) {(B) 相性（ソース）\\平均との差};
  \foreach \e/\x in {N/5.9,H/6.88,S/7.86,A/8.84}{ \node[font=\tiny] at (\x,1.75) {\e}; }
  \foreach \y/\vals in {1.3/{$-$0.10,$+$0.30,$-$0.20,0.00},0.7/{$+$0.04,$-$0.26,$+$0.19,$+$0.04},0.1/{$+$0.05,$-$0.05,$-$0.05,$+$0.05}}{
    \foreach \v [count=\j from 0] in \vals { \node[c,fill=cSrc!45,minimum width=9.4mm] at (5.9+0.98*\j,\y) {\v}; }
  }
  % (C) 水準（K 本）
  \node[font=\scriptsize\bfseries,align=center] at (10.2,2.35) {(C) 水準\\（K 本．7.2 節）};
  \foreach \v/\y in {0.617/1.3,0.507/0.7,0.477/0.1}{ \node[c,fill=cK!55,minimum width=9.5mm] at (10.2,\y) {\v}; }
    % (D) 結果
  \node[font=\scriptsize\bfseries,align=center] at (12.85,2.35) {(D) できた表\\＝ (C) ＋ (B)};
  \foreach \e/\x in {N/11.75,H/12.5,S/13.25,A/14.0}{ \node[font=\tiny] at (\x,1.75) {\e}; }
  \foreach \m/\c/\y/\vals in {顔/cV/1.3/{0.52,0.92,0.42,0.62},声/cA/0.7/{0.55,0.25,0.70,0.55},言葉/cT/0.1/{0.53,0.43,0.43,0.53}}{
    \foreach \v [count=\j from 0] in \vals { \node[c,fill=cRel!55] at (11.75+0.75*\j,\y) {\v}; }
  }
  \draw[arr] (10.75,0.7) -- (11.3,0.7);
  % 計算の例
  \node[draw,rounded corners=2pt,fill=black!3,font=\scriptsize,align=left,anchor=north west] at (-0.6,-0.85) {
    計算の例（顔の Happy）：K 本で測った顔の水準 0.617 ＋ ソースで測った顔の Happy の相性 $+$0.30 ＝ 0.917 → 0.92\\
    計算の例（声の Happy）：K 本で測った声の水準 0.507 ＋ ソースで測った声の Happy の相性 $-$0.26 ＝ 0.247 → 0.25};
\end{tikzpicture}\end{adjustbox}
\caption{表の組み立て（例の数）．(A) のソースで見込んだ表は，顔・声・言葉の水準がほぼ同じ（0.60，0.56，0.50）である．ソースの水準は捨て，K 本で測った水準 (C) に置き換える．相性 (B) はそのまま残す．できた表 (D) は，「顔が全体に強く（K 本から），特に Happy に強い（ソースから）」を表す．値が負になれば 0 にする．相性を使わない設定（β ＝ 0）では (B) を全部 0 にするので，各行は水準だけの同じ値が並ぶ．}
\end{figure}

\begin{itemize}
  \item 最後に，表の 12 個の平均で割って\textbf{重み}にする．重みが全部 1 なら積と同じである．8 節のまとめ方は，この重みを使う．
  \item \textbf{切片}：感情ごとに違う重みを掛けると，答えが特定の感情に片寄りやすい．そこで，ラベルを使わずに，ラベル無しの発話の上で確率の平均が積と同じになるように，感情ごとの足し算の数（切片）をずらす．9 節より後の結果は，この切片をずらした答えの値である．
\end{itemize}

\begin{table}[H]
\centering
\small
\caption{9 節で並べる表の測り方．}
\begin{tabular}{>{\raggedright}p{4.6cm}>{\raggedright\arraybackslash}p{10.4cm}}
\toprule
呼び名 & 表の作り方 \\
\midrule
提案（水準は K 本＋相性はソース） & 7.1〜7.3 節のとおり．n0・戻し先・β は 6 節で選ぶ \\
K 本だけで測った表 & 12 個すべてを 7.1 節の手順で K 本から測る（戻さない．ソースの相性も使わない） \\
ソースだけで測った表 & 5 節の表をそのまま使う（ターゲットのラベルを表に使わない） \\
全ラベルで測った表（参照） & 仮のターゲットのラベル無しの発話の本当のラベル（K 本以外の全部）で 12 個を測る．実際には使えない \\
\bottomrule
\end{tabular}
\end{table}

%======================================================================
\section{④ テスト：発話ごとに答えをまとめる}

2 つの分類器（顔と声）と 2 つの感情（Happy と Sad）の小さい例で，計算の仕方を示す．実際は確率の対数で計算するが，例では重みを掛けて足す形で示す．

\begin{figure}[H]
\centering
\begin{adjustbox}{scale=1.22,max width=\linewidth}\begin{tikzpicture}
  % 入力の表
  \node[font=\small\bfseries] at (1.6,1.5) {この発話への確率};
  \foreach \e/\x in {Happy/1.5,Sad/2.7}{ \node[font=\scriptsize] at (\x,0.95) {\e}; }
  \node[font=\small,cV,anchor=east] at (0.75,0.35) {顔};
  \node[font=\small,cA,anchor=east] at (0.75,-0.4) {声};
  \foreach \v/\x/\y in {0.6/1.5/0.35,0.4/2.7/0.35,0.3/1.5/-0.4,0.7/2.7/-0.4}{
    \node[draw,fill=cClf!12,minimum width=11.4mm,minimum height=6.9mm,inner sep=0pt,font=\small] at (\x,\y) {\v};
  }
  \node[font=\small\bfseries] at (6.2,1.5) {信頼度の表};
  \foreach \e/\x in {Happy/6.1,Sad/7.3}{ \node[font=\scriptsize] at (\x,0.95) {\e}; }
  \node[font=\small,cV,anchor=east] at (5.35,0.35) {顔};
  \node[font=\small,cA,anchor=east] at (5.35,-0.4) {声};
  \foreach \v/\x/\y/\d in {0.8/6.1/0.35/90,0.4/7.3/0.35/35,0.2/6.1/-0.4/12,0.6/7.3/-0.4/60}{
    \node[draw,fill=cRel!\d,minimum width=11.4mm,minimum height=6.9mm,inner sep=0pt,font=\small] at (\x,\y) {\v};
  }
  % 計算
  \node[box,fill=cV!8,align=left,minimum width=70mm,anchor=west] (h) at (-0.4,-2.0) {\textbf{「本当は Happy」と仮定}：Happy の列の信頼度を重みにする\\$0.8 \times 0.6 + 0.2 \times 0.3 = \mathbf{0.54}$};
  \node[box,fill=cA!8,align=left,minimum width=70mm,anchor=west] (s) at (-0.4,-3.4) {\textbf{「本当は Sad」と仮定}：Sad の列の信頼度を重みにする\\$0.4 \times 0.4 + 0.6 \times 0.7 = \mathbf{0.58}$};
  \node[box,fill=cTest!8,minimum width=30mm,minimum height=12mm] (ans) at (10.6,-2.7) {点の高い Sad を\\答えにする};
  \draw[arr] (h.east) -- (ans.west); \draw[arr] (s.east) -- (ans.west);
\end{tikzpicture}\end{adjustbox}
\caption{感情を順に仮定して比べる．Happy の点は「Happy に強い分類器」を重く数えて出し，Sad の点は「Sad に強い分類器」を重く数えて出す．}
\end{figure}

\noindent\textbf{実際の計算（3 つの分類器，4 つの感情の例）}：実際は確率そのものではなく，各分類器の「感情の点」（確率の対数から，4 感情の平均を引いた値）に重みを掛けて足す．重みは 7.3 節の表 (D) を 12 個の平均 0.54 で割ったものである．
\begin{center}
\footnotesize
\setlength{\tabcolsep}{5pt}
\begin{tabular}{l cccc c cccc}
\toprule
 & \multicolumn{4}{c}{この発話の確率} & & \multicolumn{4}{c}{感情の点（対数 − 平均）} \\
\cmidrule{2-5}\cmidrule{7-10}
分類器 & N & H & S & A & & N & H & S & A \\
\midrule
顔 & 0.1 & 0.6 & 0.2 & 0.1 & & $-0.62$ & $+1.17$ & $+0.07$ & $-0.62$ \\
声 & 0.2 & 0.1 & 0.5 & 0.2 & & $-0.06$ & $-0.75$ & $+0.86$ & $-0.06$ \\
言葉 & 0.3 & 0.3 & 0.2 & 0.2 & & $+0.20$ & $+0.20$ & $-0.20$ & $-0.20$ \\
\midrule
\multicolumn{6}{l}{\textbf{積}：重みは全部 1．点を足す（確率を掛けるのと同じ順になる）} & $-0.47$ & $+0.62$ & $\mathbf{+0.73}$ & $-0.88$ \\
\midrule
\multicolumn{6}{l}{重み（顔）：0.97，1.71，0.78，1.15} & & & & \\
\multicolumn{6}{l}{重み（声）：1.02，0.47，1.30，1.02} & & & & \\
\multicolumn{6}{l}{重み（言葉）：0.99，0.80，0.80，0.99} & & & & \\
\multicolumn{6}{l}{\textbf{提案の手法}：感情ごとに，その列の重みを掛けて足す} & $-0.46$ & $\mathbf{+1.82}$ & $+1.01$ & $-0.97$ \\
\bottomrule
\end{tabular}
\end{center}
積では Sad（$+0.73$）が答えになる．提案の手法では，Happy の点は「Happy に強い顔」を重く（1.71）数え，「Happy に弱い声」を軽く（0.47）数えるので，Happy（$+1.82$）が答えになる．最後に，7.3 節の切片のずらしをしてから答えを決める．


%======================================================================
\section{下調べの結果（4 コーパス）}

\textbf{下調べ}は，評価に使うコーパスのテストを開かずに，手法の形と設定を決めるための計算である．
4 つのコーパスのうち 1 つを触らずに残し，残り 3 つのうち 1 つを仮のターゲット，2 つをソースにする（5.1 節）．
成績は仮のターゲットのテストで測る．コーパスの役の割り当て 12 通り × 話者の分け方 5 通り × K 本の引き方 2 通りで，計 120 回計算した．

並べるのは，信頼度の表の測り方 3 通り（7.3 節の表）と，参照である．どれも，積と同じ 3 つの分類器の答えをまとめる．
\begin{itemize}
  \item \textbf{提案}：水準は K 本，相性はソース．
  \item \textbf{K 本だけで測った表}：12 個すべてを K 本で測る．
  \item \textbf{ソースだけで測った表}：ターゲットのラベルを表に使わない．
  \item \textbf{全ラベルで測った表（参照）}：仮のターゲットの全ラベルで測る．実際には使えない．表を正しく測れたときの上限の目安である．
\end{itemize}
\textbf{基準}：区間の下端が 0 より上で，差が 1 点以上のとき，「上回った」と数える．

\begin{figure}[H]\centering
\begin{adjustbox}{max width=\linewidth}
\begin{tikzpicture}[x=1cm,y=0.42cm]
\draw[black!12] (0.00,-4) -- (4.85,-4);
\node[font=\small,anchor=east] at (-0.05,-4) {$-4$};
\draw[black!12] (0.00,-2) -- (4.85,-2);
\node[font=\small,anchor=east] at (-0.05,-2) {$-2$};
\draw[black!12] (0.00,0) -- (4.85,0);
\node[font=\small,anchor=east] at (-0.05,0) {0};
\draw[black!12] (0.00,2) -- (4.85,2);
\node[font=\small,anchor=east] at (-0.05,2) {$+2$};
\draw[black!12] (0.00,4) -- (4.85,4);
\node[font=\small,anchor=east] at (-0.05,4) {$+4$};
\draw[black!12] (0.00,6) -- (4.85,6);
\node[font=\small,anchor=east] at (-0.05,6) {$+6$};
\draw[black!45,dashed] (0.00,1) -- (4.85,1);
\draw[black,line width=0.8pt] (0.00,0) -- (4.85,0);
\draw (0.00,-4) -- (0.00,6);
\draw[cBar1,line width=1.3pt] (0.35,1.60) -- (0.35,3.79);
\fill[cBar1] (0.35,2.70) circle (2.1pt);
\draw[cBar2,line width=1.3pt] (0.62,1.55) -- (0.62,5.35);
\fill[cBar2] (0.62,3.45) circle (2.1pt);
\draw[cBar3,line width=1.3pt] (0.89,-1.03) -- (0.89,1.70);
\fill[cBar3] (0.89,0.33) circle (2.1pt);
\node[font=\small,rotate=35,anchor=north east,inner sep=1pt] at (0.84,-4.35) {AFEW};
\draw[cBar1,line width=1.3pt] (1.45,-0.11) -- (1.45,1.21);
\fill[cBar1] (1.45,0.55) circle (2.1pt);
\draw[cBar2,line width=1.3pt] (1.72,-0.59) -- (1.72,2.70);
\fill[cBar2] (1.72,1.06) circle (2.1pt);
\draw[cBar3,line width=1.3pt] (1.99,-2.54) -- (1.99,-0.37);
\fill[cBar3] (1.99,-1.45) circle (2.1pt);
\node[font=\small,rotate=35,anchor=north east,inner sep=1pt] at (1.94,-4.35) {IEMOCAP};
\draw[cBar1,line width=1.3pt] (2.55,-0.38) -- (2.55,1.02);
\fill[cBar1] (2.55,0.32) circle (2.1pt);
\draw[cBar2,line width=1.3pt] (2.82,-1.57) -- (2.82,0.20);
\fill[cBar2] (2.82,-0.69) circle (2.1pt);
\draw[cBar3,line width=1.3pt] (3.09,0.07) -- (3.09,0.77);
\fill[cBar3] (3.09,0.42) circle (2.1pt);
\node[font=\small,rotate=35,anchor=north east,inner sep=1pt] at (3.04,-4.35) {MOSEI};
\draw[cBar1,line width=1.3pt] (3.65,-0.66) -- (3.65,0.69);
\fill[cBar1] (3.65,0.01) circle (2.1pt);
\draw[cBar2,line width=1.3pt] (3.92,-1.63) -- (3.92,0.83);
\fill[cBar2] (3.92,-0.40) circle (2.1pt);
\draw[cBar3,line width=1.3pt] (4.19,-3.22) -- (4.19,-0.69);
\fill[cBar3] (4.19,-1.95) circle (2.1pt);
\node[font=\small,rotate=35,anchor=north east,inner sep=1pt] at (4.14,-4.35) {MELD};
\node[font=\small\bfseries] at (2.43,7.10) {K = 5};
\draw[black!12] (5.80,-4) -- (10.65,-4);
\node[font=\small,anchor=east] at (5.75,-4) {$-4$};
\draw[black!12] (5.80,-2) -- (10.65,-2);
\node[font=\small,anchor=east] at (5.75,-2) {$-2$};
\draw[black!12] (5.80,0) -- (10.65,0);
\node[font=\small,anchor=east] at (5.75,0) {0};
\draw[black!12] (5.80,2) -- (10.65,2);
\node[font=\small,anchor=east] at (5.75,2) {$+2$};
\draw[black!12] (5.80,4) -- (10.65,4);
\node[font=\small,anchor=east] at (5.75,4) {$+4$};
\draw[black!12] (5.80,6) -- (10.65,6);
\node[font=\small,anchor=east] at (5.75,6) {$+6$};
\draw[black!45,dashed] (5.80,1) -- (10.65,1);
\draw[black,line width=0.8pt] (5.80,0) -- (10.65,0);
\draw (5.80,-4) -- (5.80,6);
\draw[cBar1,line width=1.3pt] (6.15,1.09) -- (6.15,4.21);
\fill[cBar1] (6.15,2.65) circle (2.1pt);
\draw[cBar2,line width=1.3pt] (6.42,0.58) -- (6.42,4.57);
\fill[cBar2] (6.42,2.57) circle (2.1pt);
\draw[cBar3,line width=1.3pt] (6.69,-0.73) -- (6.69,1.33);
\fill[cBar3] (6.69,0.30) circle (2.1pt);
\node[font=\small,rotate=35,anchor=north east,inner sep=1pt] at (6.64,-4.35) {AFEW};
\draw[cBar1,line width=1.3pt] (7.25,-0.20) -- (7.25,1.67);
\fill[cBar1] (7.25,0.74) circle (2.1pt);
\draw[cBar2,line width=1.3pt] (7.52,-1.08) -- (7.52,3.16);
\fill[cBar2] (7.52,1.04) circle (2.1pt);
\draw[cBar3,line width=1.3pt] (7.79,-2.05) -- (7.79,-1.20);
\fill[cBar3] (7.79,-1.63) circle (2.1pt);
\node[font=\small,rotate=35,anchor=north east,inner sep=1pt] at (7.74,-4.35) {IEMOCAP};
\draw[cBar1,line width=1.3pt] (8.35,-0.28) -- (8.35,0.86);
\fill[cBar1] (8.35,0.29) circle (2.1pt);
\draw[cBar2,line width=1.3pt] (8.62,-1.16) -- (8.62,0.76);
\fill[cBar2] (8.62,-0.20) circle (2.1pt);
\draw[cBar3,line width=1.3pt] (8.89,-0.03) -- (8.89,0.65);
\fill[cBar3] (8.89,0.31) circle (2.1pt);
\node[font=\small,rotate=35,anchor=north east,inner sep=1pt] at (8.84,-4.35) {MOSEI};
\draw[cBar1,line width=1.3pt] (9.45,-0.29) -- (9.45,1.46);
\fill[cBar1] (9.45,0.58) circle (2.1pt);
\draw[cBar2,line width=1.3pt] (9.72,-0.50) -- (9.72,0.90);
\fill[cBar2] (9.72,0.20) circle (2.1pt);
\draw[cBar3,line width=1.3pt] (9.99,-3.23) -- (9.99,-0.56);
\fill[cBar3] (9.99,-1.90) circle (2.1pt);
\node[font=\small,rotate=35,anchor=north east,inner sep=1pt] at (9.94,-4.35) {MELD};
\node[font=\small\bfseries] at (8.23,7.10) {K = 10};
\draw[black!12] (11.60,-4) -- (16.45,-4);
\node[font=\small,anchor=east] at (11.55,-4) {$-4$};
\draw[black!12] (11.60,-2) -- (16.45,-2);
\node[font=\small,anchor=east] at (11.55,-2) {$-2$};
\draw[black!12] (11.60,0) -- (16.45,0);
\node[font=\small,anchor=east] at (11.55,0) {0};
\draw[black!12] (11.60,2) -- (16.45,2);
\node[font=\small,anchor=east] at (11.55,2) {$+2$};
\draw[black!12] (11.60,4) -- (16.45,4);
\node[font=\small,anchor=east] at (11.55,4) {$+4$};
\draw[black!12] (11.60,6) -- (16.45,6);
\node[font=\small,anchor=east] at (11.55,6) {$+6$};
\draw[black!45,dashed] (11.60,1) -- (16.45,1);
\draw[black,line width=0.8pt] (11.60,0) -- (16.45,0);
\draw (11.60,-4) -- (11.60,6);
\draw[cBar1,line width=1.3pt] (11.95,1.06) -- (11.95,3.70);
\fill[cBar1] (11.95,2.38) circle (2.1pt);
\draw[cBar2,line width=1.3pt] (12.22,1.63) -- (12.22,4.60);
\fill[cBar2] (12.22,3.11) circle (2.1pt);
\draw[cBar3,line width=1.3pt] (12.49,-0.76) -- (12.49,1.45);
\fill[cBar3] (12.49,0.35) circle (2.1pt);
\node[font=\small,rotate=35,anchor=north east,inner sep=1pt] at (12.44,-4.35) {AFEW};
\draw[cBar1,line width=1.3pt] (13.05,-0.00) -- (13.05,0.75);
\fill[cBar1] (13.05,0.37) circle (2.1pt);
\draw[cBar2,line width=1.3pt] (13.32,-0.33) -- (13.32,1.89);
\fill[cBar2] (13.32,0.78) circle (2.1pt);
\draw[cBar3,line width=1.3pt] (13.59,-2.61) -- (13.59,-1.06);
\fill[cBar3] (13.59,-1.84) circle (2.1pt);
\node[font=\small,rotate=35,anchor=north east,inner sep=1pt] at (13.54,-4.35) {IEMOCAP};
\draw[cBar1,line width=1.3pt] (14.15,-0.19) -- (14.15,0.79);
\fill[cBar1] (14.15,0.30) circle (2.1pt);
\draw[cBar2,line width=1.3pt] (14.42,-0.51) -- (14.42,0.85);
\fill[cBar2] (14.42,0.17) circle (2.1pt);
\draw[cBar3,line width=1.3pt] (14.69,-0.16) -- (14.69,0.45);
\fill[cBar3] (14.69,0.14) circle (2.1pt);
\node[font=\small,rotate=35,anchor=north east,inner sep=1pt] at (14.64,-4.35) {MOSEI};
\draw[cBar1,line width=1.3pt] (15.25,0.08) -- (15.25,1.97);
\fill[cBar1] (15.25,1.03) circle (2.1pt);
\draw[cBar2,line width=1.3pt] (15.52,0.06) -- (15.52,1.33);
\fill[cBar2] (15.52,0.69) circle (2.1pt);
\draw[cBar3,line width=1.3pt] (15.79,-3.21) -- (15.79,-1.51);
\fill[cBar3] (15.79,-2.36) circle (2.1pt);
\node[font=\small,rotate=35,anchor=north east,inner sep=1pt] at (15.74,-4.35) {MELD};
\node[font=\small\bfseries] at (14.03,7.10) {K = 25};
\node[font=\small,rotate=90] at (-0.95,1.00) {積との差（UAR の点）};
\fill[cBar1] (0.12,-9.00) circle (2.4pt);
\draw[cBar1,line width=1.3pt] (0.12,-9.80) -- (0.12,-8.20);
\node[font=\small,anchor=west] at (0.30,-9.00) {提案（水準は K 本＋相性はソース）};
\fill[cBar2] (7.50,-9.00) circle (2.4pt);
\draw[cBar2,line width=1.3pt] (7.50,-9.80) -- (7.50,-8.20);
\node[font=\small,anchor=west] at (7.68,-9.00) {K 本だけで測った表};
\fill[cBar3] (12.00,-9.00) circle (2.4pt);
\draw[cBar3,line width=1.3pt] (12.00,-9.80) -- (12.00,-8.20);
\node[font=\small,anchor=west] at (12.18,-9.00) {ソースだけで測った表};
\node[font=\small,anchor=west] at (0,-11.20) {点は差，縦の線は 95\% の区間．太い横線は差 0，点線は $+1$ 点（基準）．};
\end{tikzpicture}
\end{adjustbox}
\caption{下調べ（4 コーパス）．表の測り方 3 通りの，積との差．}
\end{figure}

\begin{table}[H]\centering\small
\caption{下調べ（4 コーパス）：積との差．括弧は 95\% の区間．太字と $\uparrow$ は基準（区間の下端が 0 より上で，差が 1 点以上）を満たした条件．}
\setlength{\tabcolsep}{5pt}
\begin{adjustbox}{max width=\linewidth}
\begin{tabular}{lllll}
\toprule
K & AFEW & IEMOCAP & MOSEI & MELD \\
\midrule
\multicolumn{5}{l}{\textbf{提案 − 積}} \\
5 & \textbf{$+2.70$} {\scriptsize[$+1.6$, $+3.8$]} $\uparrow$ & $+0.55$ {\scriptsize[$-0.1$, $+1.2$]} & $+0.32$ {\scriptsize[$-0.4$, $+1.0$]} & $+0.01$ {\scriptsize[$-0.7$, $+0.7$]} \\
10 & \textbf{$+2.65$} {\scriptsize[$+1.1$, $+4.2$]} $\uparrow$ & $+0.74$ {\scriptsize[$-0.2$, $+1.7$]} & $+0.29$ {\scriptsize[$-0.3$, $+0.9$]} & $+0.58$ {\scriptsize[$-0.3$, $+1.5$]} \\
25 & \textbf{$+2.38$} {\scriptsize[$+1.1$, $+3.7$]} $\uparrow$ & $+0.37$ {\scriptsize[$-0.0$, $+0.8$]} & $+0.30$ {\scriptsize[$-0.2$, $+0.8$]} & \textbf{$+1.03$} {\scriptsize[$+0.1$, $+2.0$]} $\uparrow$ \\
\midrule
\multicolumn{5}{l}{\textbf{全ラベルで測った表 − 積（参照）}} \\
5 & \textbf{$+4.75$} {\scriptsize[$+2.6$, $+6.9$]} $\uparrow$ & \textbf{$+1.25$} {\scriptsize[$+0.1$, $+2.4$]} $\uparrow$ & \textbf{$+1.59$} {\scriptsize[$+1.0$, $+2.2$]} $\uparrow$ & \textbf{$+1.11$} {\scriptsize[$+0.1$, $+2.1$]} $\uparrow$ \\
10 & \textbf{$+3.96$} {\scriptsize[$+1.8$, $+6.1$]} $\uparrow$ & $+1.10$ {\scriptsize[$-0.4$, $+2.6$]} & \textbf{$+1.11$} {\scriptsize[$+0.7$, $+1.5$]} $\uparrow$ & $+0.93$ {\scriptsize[$+0.0$, $+1.9$]} \\
25 & \textbf{$+3.93$} {\scriptsize[$+1.6$, $+6.3$]} $\uparrow$ & \textbf{$+1.19$} {\scriptsize[$+0.3$, $+2.1$]} $\uparrow$ & $+0.63$ {\scriptsize[$+0.3$, $+0.9$]} & \textbf{$+1.23$} {\scriptsize[$+0.6$, $+1.8$]} $\uparrow$ \\
\bottomrule
\end{tabular}
\end{adjustbox}
\end{table}

\begin{table}[H]\centering\small
\caption{下調べ（4 コーパス）：UAR（点）．}
\begin{tabular}{llcccc}
\toprule
K & 答えの作り方 & AFEW & IEMOCAP & MOSEI & MELD \\
\midrule
5 & 積 & 68.6 & 71.2 & 40.7 & 50.3 \\
 & 提案 & 71.3 & 71.8 & 41.0 & 50.3 \\
 & 全ラベルで測った表（参照） & 73.3 & 72.5 & 42.3 & 51.4 \\
 & K 本だけで学んだ融合の分類器 & 69.4 & 70.4 & 40.9 & 47.4 \\
\midrule
10 & 積 & 71.0 & 72.5 & 41.6 & 50.3 \\
 & 提案 & 73.6 & 73.2 & 41.8 & 50.9 \\
 & 全ラベルで測った表（参照） & 74.9 & 73.6 & 42.7 & 51.2 \\
 & K 本だけで学んだ融合の分類器 & 71.8 & 72.0 & 41.3 & 48.0 \\
\midrule
25 & 積 & 72.2 & 75.7 & 43.5 & 51.7 \\
 & 提案 & 74.6 & 76.1 & 43.8 & 52.7 \\
 & 全ラベルで測った表（参照） & 76.1 & 76.9 & 44.1 & 52.9 \\
 & K 本だけで学んだ融合の分類器 & 74.2 & 75.4 & 43.2 & 49.8 \\
\bottomrule
\end{tabular}
\end{table}


%======================================================================
\section{コーパスを 6 つに増やした結果}

CREMA-D（声と顔）と RAVDESS（声だけ）を足して，同じ下調べを 6 コーパスで流した（300 回）．どちらも，俳優が決まった文を演技で読むコーパスである．
\begin{itemize}
  \item 元の 4 コーパスの値も，9 節の値と同じにはならない．ソースに CREMA-D と RAVDESS が入り，出発点と微調整した分類器が変わるためである．
  \item RAVDESS は声しか無いので，3 つの答えをまとめる提案は使えない．図と表には入れていない．
  \item この計算では，提案の答えを疑似ラベルにして融合の分類器を学び直すステップも流した（最後の表）．
  \item 出力の欠けと重複の点検は通った．結果の文の照合はこれからである．
\end{itemize}

\begin{figure}[H]\centering
\begin{adjustbox}{max width=\linewidth}
\begin{tikzpicture}[x=1cm,y=0.42cm]
\draw[black!12] (0.00,-4) -- (5.05,-4);
\node[font=\small,anchor=east] at (-0.05,-4) {$-4$};
\draw[black!12] (0.00,-2) -- (5.05,-2);
\node[font=\small,anchor=east] at (-0.05,-2) {$-2$};
\draw[black!12] (0.00,0) -- (5.05,0);
\node[font=\small,anchor=east] at (-0.05,0) {0};
\draw[black!12] (0.00,2) -- (5.05,2);
\node[font=\small,anchor=east] at (-0.05,2) {$+2$};
\draw[black!12] (0.00,4) -- (5.05,4);
\node[font=\small,anchor=east] at (-0.05,4) {$+4$};
\draw[black!12] (0.00,6) -- (5.05,6);
\node[font=\small,anchor=east] at (-0.05,6) {$+6$};
\draw[black!45,dashed] (0.00,1) -- (5.05,1);
\draw[black,line width=0.8pt] (0.00,0) -- (5.05,0);
\draw (0.00,-4) -- (0.00,6);
\draw[cBar1,line width=1.3pt] (0.35,1.18) -- (0.35,3.16);
\fill[cBar1] (0.35,2.17) circle (2.1pt);
\draw[cBar2,line width=1.3pt] (0.62,1.63) -- (0.62,4.91);
\fill[cBar2] (0.62,3.27) circle (2.1pt);
\draw[cBar3,line width=1.3pt] (0.89,0.34) -- (0.89,2.20);
\fill[cBar3] (0.89,1.27) circle (2.1pt);
\node[font=\small,rotate=35,anchor=north east,inner sep=1pt] at (0.84,-4.35) {AFEW};
\draw[cBar1,line width=1.3pt] (1.27,-0.66) -- (1.27,1.07);
\fill[cBar1] (1.27,0.21) circle (2.1pt);
\draw[cBar2,line width=1.3pt] (1.54,-0.52) -- (1.54,3.01);
\fill[cBar2] (1.54,1.24) circle (2.1pt);
\draw[cBar3,line width=1.3pt] (1.81,-1.68) -- (1.81,0.67);
\fill[cBar3] (1.81,-0.51) circle (2.1pt);
\node[font=\small,rotate=35,anchor=north east,inner sep=1pt] at (1.76,-4.35) {IEMOCAP};
\draw[cBar1,line width=1.3pt] (2.19,-0.46) -- (2.19,0.64);
\fill[cBar1] (2.19,0.09) circle (2.1pt);
\draw[cBar2,line width=1.3pt] (2.46,-1.79) -- (2.46,-0.15);
\fill[cBar2] (2.46,-0.97) circle (2.1pt);
\draw[cBar3,line width=1.3pt] (2.73,0.07) -- (2.73,0.65);
\fill[cBar3] (2.73,0.36) circle (2.1pt);
\node[font=\small,rotate=35,anchor=north east,inner sep=1pt] at (2.68,-4.35) {MOSEI};
\draw[cBar1,line width=1.3pt] (3.11,-0.82) -- (3.11,2.08);
\fill[cBar1] (3.11,0.63) circle (2.1pt);
\draw[cBar2,line width=1.3pt] (3.38,-1.76) -- (3.38,1.56);
\fill[cBar2] (3.38,-0.10) circle (2.1pt);
\draw[cBar3,line width=1.3pt] (3.65,-3.04) -- (3.65,-0.75);
\fill[cBar3] (3.65,-1.90) circle (2.1pt);
\node[font=\small,rotate=35,anchor=north east,inner sep=1pt] at (3.60,-4.35) {MELD};
\draw[cBar1,line width=1.3pt] (4.03,0.05) -- (4.03,0.71);
\fill[cBar1] (4.03,0.38) circle (2.1pt);
\draw[cBar2,line width=1.3pt] (4.30,-0.21) -- (4.30,0.68);
\fill[cBar2] (4.30,0.23) circle (2.1pt);
\draw[cBar3,line width=1.3pt] (4.57,-0.02) -- (4.57,0.59);
\fill[cBar3] (4.57,0.28) circle (2.1pt);
\node[font=\small,rotate=35,anchor=north east,inner sep=1pt] at (4.52,-4.35) {CREMA-D};
\node[font=\small\bfseries] at (2.53,7.10) {K = 5};
\draw[black!12] (6.00,-4) -- (11.05,-4);
\node[font=\small,anchor=east] at (5.95,-4) {$-4$};
\draw[black!12] (6.00,-2) -- (11.05,-2);
\node[font=\small,anchor=east] at (5.95,-2) {$-2$};
\draw[black!12] (6.00,0) -- (11.05,0);
\node[font=\small,anchor=east] at (5.95,0) {0};
\draw[black!12] (6.00,2) -- (11.05,2);
\node[font=\small,anchor=east] at (5.95,2) {$+2$};
\draw[black!12] (6.00,4) -- (11.05,4);
\node[font=\small,anchor=east] at (5.95,4) {$+4$};
\draw[black!12] (6.00,6) -- (11.05,6);
\node[font=\small,anchor=east] at (5.95,6) {$+6$};
\draw[black!45,dashed] (6.00,1) -- (11.05,1);
\draw[black,line width=0.8pt] (6.00,0) -- (11.05,0);
\draw (6.00,-4) -- (6.00,6);
\draw[cBar1,line width=1.3pt] (6.35,1.44) -- (6.35,4.41);
\fill[cBar1] (6.35,2.92) circle (2.1pt);
\draw[cBar2,line width=1.3pt] (6.62,1.36) -- (6.62,4.43);
\fill[cBar2] (6.62,2.90) circle (2.1pt);
\draw[cBar3,line width=1.3pt] (6.89,-0.16) -- (6.89,2.41);
\fill[cBar3] (6.89,1.12) circle (2.1pt);
\node[font=\small,rotate=35,anchor=north east,inner sep=1pt] at (6.84,-4.35) {AFEW};
\draw[cBar1,line width=1.3pt] (7.27,-0.37) -- (7.27,1.36);
\fill[cBar1] (7.27,0.50) circle (2.1pt);
\draw[cBar2,line width=1.3pt] (7.54,-0.26) -- (7.54,1.88);
\fill[cBar2] (7.54,0.81) circle (2.1pt);
\draw[cBar3,line width=1.3pt] (7.81,-1.74) -- (7.81,1.00);
\fill[cBar3] (7.81,-0.37) circle (2.1pt);
\node[font=\small,rotate=35,anchor=north east,inner sep=1pt] at (7.76,-4.35) {IEMOCAP};
\draw[cBar1,line width=1.3pt] (8.19,-0.27) -- (8.19,0.87);
\fill[cBar1] (8.19,0.30) circle (2.1pt);
\draw[cBar2,line width=1.3pt] (8.46,-0.99) -- (8.46,0.57);
\fill[cBar2] (8.46,-0.21) circle (2.1pt);
\draw[cBar3,line width=1.3pt] (8.73,0.09) -- (8.73,0.58);
\fill[cBar3] (8.73,0.33) circle (2.1pt);
\node[font=\small,rotate=35,anchor=north east,inner sep=1pt] at (8.68,-4.35) {MOSEI};
\draw[cBar1,line width=1.3pt] (9.11,-0.63) -- (9.11,0.92);
\fill[cBar1] (9.11,0.15) circle (2.1pt);
\draw[cBar2,line width=1.3pt] (9.38,-1.03) -- (9.38,1.48);
\fill[cBar2] (9.38,0.23) circle (2.1pt);
\draw[cBar3,line width=1.3pt] (9.65,-2.89) -- (9.65,-1.30);
\fill[cBar3] (9.65,-2.09) circle (2.1pt);
\node[font=\small,rotate=35,anchor=north east,inner sep=1pt] at (9.60,-4.35) {MELD};
\draw[cBar1,line width=1.3pt] (10.03,0.26) -- (10.03,0.86);
\fill[cBar1] (10.03,0.56) circle (2.1pt);
\draw[cBar2,line width=1.3pt] (10.30,0.20) -- (10.30,1.12);
\fill[cBar2] (10.30,0.66) circle (2.1pt);
\draw[cBar3,line width=1.3pt] (10.57,-0.04) -- (10.57,0.48);
\fill[cBar3] (10.57,0.22) circle (2.1pt);
\node[font=\small,rotate=35,anchor=north east,inner sep=1pt] at (10.52,-4.35) {CREMA-D};
\node[font=\small\bfseries] at (8.53,7.10) {K = 10};
\draw[black!12] (12.00,-4) -- (17.05,-4);
\node[font=\small,anchor=east] at (11.95,-4) {$-4$};
\draw[black!12] (12.00,-2) -- (17.05,-2);
\node[font=\small,anchor=east] at (11.95,-2) {$-2$};
\draw[black!12] (12.00,0) -- (17.05,0);
\node[font=\small,anchor=east] at (11.95,0) {0};
\draw[black!12] (12.00,2) -- (17.05,2);
\node[font=\small,anchor=east] at (11.95,2) {$+2$};
\draw[black!12] (12.00,4) -- (17.05,4);
\node[font=\small,anchor=east] at (11.95,4) {$+4$};
\draw[black!12] (12.00,6) -- (17.05,6);
\node[font=\small,anchor=east] at (11.95,6) {$+6$};
\draw[black!45,dashed] (12.00,1) -- (17.05,1);
\draw[black,line width=0.8pt] (12.00,0) -- (17.05,0);
\draw (12.00,-4) -- (12.00,6);
\draw[cBar1,line width=1.3pt] (12.35,1.60) -- (12.35,3.83);
\fill[cBar1] (12.35,2.72) circle (2.1pt);
\draw[cBar2,line width=1.3pt] (12.62,1.51) -- (12.62,3.93);
\fill[cBar2] (12.62,2.72) circle (2.1pt);
\draw[cBar3,line width=1.3pt] (12.89,-0.90) -- (12.89,2.22);
\fill[cBar3] (12.89,0.66) circle (2.1pt);
\node[font=\small,rotate=35,anchor=north east,inner sep=1pt] at (12.84,-4.35) {AFEW};
\draw[cBar1,line width=1.3pt] (13.27,-0.13) -- (13.27,1.02);
\fill[cBar1] (13.27,0.44) circle (2.1pt);
\draw[cBar2,line width=1.3pt] (13.54,-0.55) -- (13.54,1.60);
\fill[cBar2] (13.54,0.52) circle (2.1pt);
\draw[cBar3,line width=1.3pt] (13.81,-1.42) -- (13.81,0.13);
\fill[cBar3] (13.81,-0.64) circle (2.1pt);
\node[font=\small,rotate=35,anchor=north east,inner sep=1pt] at (13.76,-4.35) {IEMOCAP};
\draw[cBar1,line width=1.3pt] (14.19,-0.19) -- (14.19,0.84);
\fill[cBar1] (14.19,0.33) circle (2.1pt);
\draw[cBar2,line width=1.3pt] (14.46,-0.50) -- (14.46,0.78);
\fill[cBar2] (14.46,0.14) circle (2.1pt);
\draw[cBar3,line width=1.3pt] (14.73,-0.02) -- (14.73,0.66);
\fill[cBar3] (14.73,0.32) circle (2.1pt);
\node[font=\small,rotate=35,anchor=north east,inner sep=1pt] at (14.68,-4.35) {MOSEI};
\draw[cBar1,line width=1.3pt] (15.11,0.19) -- (15.11,1.76);
\fill[cBar1] (15.11,0.97) circle (2.1pt);
\draw[cBar2,line width=1.3pt] (15.38,0.15) -- (15.38,1.45);
\fill[cBar2] (15.38,0.80) circle (2.1pt);
\draw[cBar3,line width=1.3pt] (15.65,-3.81) -- (15.65,-1.14);
\fill[cBar3] (15.65,-2.47) circle (2.1pt);
\node[font=\small,rotate=35,anchor=north east,inner sep=1pt] at (15.60,-4.35) {MELD};
\draw[cBar1,line width=1.3pt] (16.03,0.21) -- (16.03,0.91);
\fill[cBar1] (16.03,0.56) circle (2.1pt);
\draw[cBar2,line width=1.3pt] (16.30,0.06) -- (16.30,1.05);
\fill[cBar2] (16.30,0.55) circle (2.1pt);
\draw[cBar3,line width=1.3pt] (16.57,-0.19) -- (16.57,0.58);
\fill[cBar3] (16.57,0.19) circle (2.1pt);
\node[font=\small,rotate=35,anchor=north east,inner sep=1pt] at (16.52,-4.35) {CREMA-D};
\node[font=\small\bfseries] at (14.53,7.10) {K = 25};
\node[font=\small,rotate=90] at (-0.95,1.00) {積との差（UAR の点）};
\fill[cBar1] (0.12,-9.00) circle (2.4pt);
\draw[cBar1,line width=1.3pt] (0.12,-9.80) -- (0.12,-8.20);
\node[font=\small,anchor=west] at (0.30,-9.00) {提案（水準は K 本＋相性はソース）};
\fill[cBar2] (7.50,-9.00) circle (2.4pt);
\draw[cBar2,line width=1.3pt] (7.50,-9.80) -- (7.50,-8.20);
\node[font=\small,anchor=west] at (7.68,-9.00) {K 本だけで測った表};
\fill[cBar3] (12.00,-9.00) circle (2.4pt);
\draw[cBar3,line width=1.3pt] (12.00,-9.80) -- (12.00,-8.20);
\node[font=\small,anchor=west] at (12.18,-9.00) {ソースだけで測った表};
\node[font=\small,anchor=west] at (0,-11.20) {点は差，縦の線は 95\% の区間．太い横線は差 0，点線は $+1$ 点（基準）．};
\end{tikzpicture}
\end{adjustbox}
\caption{6 コーパス．表の測り方 3 通りの，積との差．RAVDESS は声しか無いので，提案を使えない．}
\end{figure}

\begin{table}[H]\centering\small
\caption{6 コーパス：積との差．印と括弧の意味は 9 節の表と同じ．}
\setlength{\tabcolsep}{5pt}
\begin{adjustbox}{max width=\linewidth}
\begin{tabular}{llllll}
\toprule
K & AFEW & IEMOCAP & MOSEI & MELD & CREMA-D \\
\midrule
\multicolumn{6}{l}{\textbf{提案 − 積}} \\
5 & \textbf{$+2.17$} {\scriptsize[$+1.2$, $+3.2$]} $\uparrow$ & $+0.21$ {\scriptsize[$-0.7$, $+1.1$]} & $+0.09$ {\scriptsize[$-0.5$, $+0.6$]} & $+0.63$ {\scriptsize[$-0.8$, $+2.1$]} & $+0.38$ {\scriptsize[$+0.1$, $+0.7$]} \\
10 & \textbf{$+2.92$} {\scriptsize[$+1.4$, $+4.4$]} $\uparrow$ & $+0.50$ {\scriptsize[$-0.4$, $+1.4$]} & $+0.30$ {\scriptsize[$-0.3$, $+0.9$]} & $+0.15$ {\scriptsize[$-0.6$, $+0.9$]} & $+0.56$ {\scriptsize[$+0.3$, $+0.9$]} \\
25 & \textbf{$+2.72$} {\scriptsize[$+1.6$, $+3.8$]} $\uparrow$ & $+0.44$ {\scriptsize[$-0.1$, $+1.0$]} & $+0.33$ {\scriptsize[$-0.2$, $+0.8$]} & $+0.97$ {\scriptsize[$+0.2$, $+1.8$]} & $+0.56$ {\scriptsize[$+0.2$, $+0.9$]} \\
\midrule
\multicolumn{6}{l}{\textbf{全ラベルで測った表 − 積（参照）}} \\
5 & \textbf{$+4.31$} {\scriptsize[$+3.0$, $+5.6$]} $\uparrow$ & \textbf{$+1.16$} {\scriptsize[$+0.1$, $+2.2$]} $\uparrow$ & \textbf{$+1.32$} {\scriptsize[$+0.7$, $+2.0$]} $\uparrow$ & \textbf{$+1.13$} {\scriptsize[$+0.1$, $+2.2$]} $\uparrow$ & $+0.79$ {\scriptsize[$+0.3$, $+1.3$]} \\
10 & \textbf{$+3.71$} {\scriptsize[$+2.0$, $+5.4$]} $\uparrow$ & \textbf{$+1.17$} {\scriptsize[$+0.2$, $+2.2$]} $\uparrow$ & $+0.99$ {\scriptsize[$+0.6$, $+1.4$]} & $+0.95$ {\scriptsize[$+0.3$, $+1.6$]} & $+0.87$ {\scriptsize[$+0.5$, $+1.3$]} \\
25 & \textbf{$+3.55$} {\scriptsize[$+1.1$, $+5.9$]} $\uparrow$ & $+0.90$ {\scriptsize[$-0.2$, $+2.0$]} & $+0.62$ {\scriptsize[$+0.3$, $+0.9$]} & \textbf{$+1.23$} {\scriptsize[$+0.6$, $+1.9$]} $\uparrow$ & $+0.71$ {\scriptsize[$+0.2$, $+1.2$]} \\
\bottomrule
\end{tabular}
\end{adjustbox}
\end{table}

\begin{table}[H]\centering\small
\caption{6 コーパス：UAR（点）．CREMA-D の融合は声と顔の 2 つ．}
\begin{tabular}{llccccc}
\toprule
K & 答えの作り方 & AFEW & IEMOCAP & MOSEI & MELD & CREMA-D \\
\midrule
5 & 積 & 68.8 & 71.8 & 41.1 & 50.5 & 85.1 \\
 & 提案 & 71.0 & 72.0 & 41.1 & 51.1 & 85.4 \\
 & 全ラベルで測った表（参照） & 73.1 & 72.9 & 42.4 & 51.7 & 85.9 \\
 & K 本だけで学んだ融合の分類器 & 69.5 & 70.7 & 41.1 & 47.8 & 82.0 \\
\midrule
10 & 積 & 70.7 & 73.4 & 41.8 & 50.8 & 85.6 \\
 & 提案 & 73.7 & 73.9 & 42.1 & 51.0 & 86.2 \\
 & 全ラベルで測った表（参照） & 74.4 & 74.6 & 42.8 & 51.7 & 86.5 \\
 & K 本だけで学んだ融合の分類器 & 71.7 & 72.5 & 41.5 & 48.5 & 82.9 \\
\midrule
25 & 積 & 72.6 & 76.3 & 43.7 & 51.9 & 87.0 \\
 & 提案 & 75.3 & 76.7 & 44.1 & 52.8 & 87.5 \\
 & 全ラベルで測った表（参照） & 76.1 & 77.2 & 44.4 & 53.1 & 87.7 \\
 & K 本だけで学んだ融合の分類器 & 74.2 & 75.7 & 43.3 & 50.1 & 84.9 \\
\bottomrule
\end{tabular}
\end{table}

\begin{table}[H]\centering\small
\caption{6 コーパス：学び直した融合の分類器の差（提案の答えを疑似ラベルにした分類器 − 積を疑似ラベルにした分類器）．}
\setlength{\tabcolsep}{5pt}
\begin{adjustbox}{max width=\linewidth}
\begin{tabular}{llllll}
\toprule
K & AFEW & IEMOCAP & MOSEI & MELD & CREMA-D \\
\midrule
5 & \textbf{$+1.81$} {\scriptsize[$+0.8$, $+2.8$]} $\uparrow$ & $+0.23$ {\scriptsize[$-0.5$, $+0.9$]} & $+0.19$ {\scriptsize[$-0.4$, $+0.7$]} & $+0.66$ {\scriptsize[$-0.5$, $+1.8$]} & $+0.49$ {\scriptsize[$+0.2$, $+0.8$]} \\
10 & \textbf{$+1.75$} {\scriptsize[$+0.6$, $+2.9$]} $\uparrow$ & $+0.24$ {\scriptsize[$-0.5$, $+1.0$]} & $+0.24$ {\scriptsize[$-0.3$, $+0.8$]} & $+0.30$ {\scriptsize[$-0.3$, $+0.9$]} & $+0.40$ {\scriptsize[$+0.1$, $+0.7$]} \\
25 & \textbf{$+1.16$} {\scriptsize[$+0.3$, $+2.0$]} $\uparrow$ & $+0.28$ {\scriptsize[$-0.5$, $+1.0$]} & $+0.31$ {\scriptsize[$-0.2$, $+0.8$]} & \textbf{$+1.14$} {\scriptsize[$+0.3$, $+2.0$]} $\uparrow$ & $+0.43$ {\scriptsize[$+0.1$, $+0.8$]} \\
\bottomrule
\end{tabular}
\end{adjustbox}
\end{table}


%======================================================================
\section{分かったことと，まだ分からないこと}

\subsection*{分かったこと（9 節と 10 節の数から）}

\begin{itemize}
\item \textbf{提案は AFEW で積を上回った．}差は，4 コーパスの下調べで $+2.38$〜$+2.70$ 点，6 コーパスで $+2.17$〜$+2.92$ 点で，基準を満たしたのは，4 コーパスの下調べで 3 つの K のうち 3 つ，6 コーパスで 3 つだった．
\item \textbf{ほかのコーパスでは差が小さい．}AFEW 以外の差は，4 コーパスの下調べで $+0.01$〜$+1.03$ 点（9 条件），6 コーパスで $+0.09$〜$+0.97$ 点（12 条件）だった．差が負だった条件は 0 個である．基準を満たしたのは，4 コーパスの下調べの MELD の K = 25 だけで，6 コーパスでは 0 条件だった．
\item CREMA-D（声と顔）の差は $+0.38$〜$+0.56$ 点で，3 つの K とも区間は 0 より上だった（3 条件）．差が 1 点に届かないので，基準は満たさない．
\item \textbf{水準を K 本で測り，相性をソースから持ち込む組み合わせは，片方だけの版の損を避けた．}ソースだけで測った表は，MELD で積を下回った（4 コーパスの下調べで $-2.36$〜$-1.90$ 点，6 コーパスで $-2.47$〜$-1.90$ 点．基準を満たして下回ったのは 6 条件）．K 本だけで測った表は，6 コーパスの MOSEI の K = 5（$-0.97$ 点，区間 $-1.8$〜$-0.1$）で区間が 0 より下だった．同じ条件での提案の差は $+0.01$〜$+1.03$ 点で，負の値は無かった．
\item \textbf{「どの感情でどのモダリティが強いか」の成分の寄与は小さい．}提案からこの成分だけを消した版との差は，4 コーパスの下調べで $-0.02$〜$+0.36$ 点，6 コーパスで $-0.15$〜$+0.36$ 点で，基準を満たした条件は 0 個だった．提案と積の差は，主にモダリティごとの重み（水準）から来ている．
\item \textbf{表を全ラベルで測れたとしたときの余地}（実際には使えない参照）：積との差は，AFEW で $+3.55$〜$+4.75$ 点，ほかのコーパスで $+0.62$〜$+1.59$ 点だった．AFEW 以外では，表を全ラベルで測っても，積との差は 1.6 点以下だった．
\item \textbf{提案の答えを疑似ラベルにして融合の分類器を学び直すと，差は AFEW で残った}（6 コーパス）．積を疑似ラベルにした分類器との差は，AFEW で $+1.16$〜$+1.81$ 点（3 つの K とも基準を満たす），ほかのコーパスで $+0.19$〜$+1.14$ 点だった（基準を満たしたのは MELD の K = 25）．
\end{itemize}


\subsection*{まだ分からないこと}
\begin{itemize}
  \item 下調べはソースが少なく（4 コーパスの下調べでは 2 つ），K 本の引き方は 2 通りである．ソースが 3 つで，引き方が 10 通りの評価の回で同じ向きになるかは，12 節の計算で測る．
  \item 既にある方法（確率の平均，温度の較正，重みの表を乱数で探す方法など）と比べてどうかは，まだ測っていない．12 節の計算で測る．
  \item 提案の形は，下調べの結果を見て選んだ．評価の回のテストの発話は，下調べで仮のターゲットにしたときのテストと同じ発話である．評価の回は，見ていないデータでの確かめにはならない．
\end{itemize}

%======================================================================
\section{いま計算している評価の回}

\begin{waku}{状態}
2026 年 11 月 14 日 21 時 25 分に研究室の PC で起動した．単位の計算が 5.3〜5.8 時間，集計が 10〜30 分の見込みである．結果が出たら，この節を結果に書き換える．
\end{waku}

\subsection*{何を測るか}
\begin{itemize}
  \item 4 つのコーパスのそれぞれをターゲットにし，残り 3 つをソースにする．話者の分け方 5 通り × K 本の引き方 10 通りで，1 コーパスあたり 50 回，計 200 回計算する．K は 5・10・25．
  \item 設定（戻す強さ n0，戻し先，相性を使うか）は，下調べが選んだ値に固定した（下の表）．評価の回では選び直さない．
  \item 主にする K は置かない．上回った・上回らないの判定の規則も置かず，結果を表で見る．
\end{itemize}
\begin{center}
\small
\begin{adjustbox}{max width=\linewidth}
\begin{tabular}{llll}
\toprule
ターゲット & K = 5 & K = 10 & K = 25 \\
\midrule
AFEW & n0 ＝ 20，戻し先は 3 つの平均 & n0 ＝ 0 & n0 ＝ 20，戻し先は 3 つの平均 \\
IEMOCAP & n0 ＝ 0 & n0 ＝ 0 & n0 ＝ 0 \\
MOSEI & n0 ＝ 0 & n0 ＝ 0 & n0 ＝ 0 \\
MELD & n0 ＝ 20，戻し先はソースの水準 & n0 ＝ 20，戻し先はソースの水準 & n0 ＝ 0 \\
\bottomrule
\end{tabular}
\end{adjustbox}

\smallskip
どの条件でも相性は使う．n0 ＝ 0 は，K 本で測った水準をそのまま使う設定である．
\end{center}

\subsection*{何と比べるか}
\begin{center}
\small
\begin{tabular}{>{\raggedright}p{3.9cm}>{\raggedright\arraybackslash}p{11.2cm}}
\toprule
組 & 相手 \\
\midrule
同じ 3 つの分類器の答えの，別のまとめ方 & 積／確率の平均／発話ごとに自信のある分類器を重くする平均／水準だけの重み／K 本だけで測った表／取り違えの表で確率を読み替えて掛ける方法／分類器ごとの温度の較正／K 本で合わせたモダリティごとの重み／重みの表を乱数で探す方法（Kaya ほか 2017 の型．候補の数は 10・100・1,000）／別のコーパスで選んだ指数の積 \\
学び直した 1 つの分類器 & K 本だけで学んだ融合の分類器／積を疑似ラベルにして学び直した融合の分類器／温度つき蒸留／一致した発話だけの疑似ラベル／ラベルと疑似ラベルを混ぜる方法（Zhang ほか 2021 の形）／共訓練／共正則化．同じ 200 回の上の差と区間を出す \\
標準の手法 & DANN，CORAL，MME，自己学習など（前に計算した結果の UAR を並べる） \\
\bottomrule
\end{tabular}
\end{center}
提案の部品を 1 つずつ外した版も測る．外す部品は，相性，分類器どうしの水準の違い，水準を戻す操作，残り方，「どの感情でどのモダリティが強いか」の成分である．差がどの部品から来たかを分けて読むために使う．

%======================================================================
\section{用語}
\begin{center}
\small
\begin{tabular}{>{\raggedright}p{3.6cm}p{11.8cm}}
\toprule
用語 & 意味 \\
\midrule
ソース／ターゲット & ラベルが全部ある既存のコーパス／ラベルが少ない新しいコーパス \\
K 本 & ターゲットで感情ごとに付いているラベルの本数（5，10 か 25） \\
分類器 & 発話の特徴から 4 つの感情の確率を出す線形の式 \\
微調整 & ソースの分類器の平均から始めて，K 本で学び直すこと \\
積 & 微調整した声・顔・言葉の確率を同じ強さで掛ける，今までのまとめ方 \\
提案の手法 & 信頼度の表から作った重みで，3 つの分類器の点を感情ごとに足すまとめ方．水準は K 本，相性はソースで測る \\
信頼度 & 分類器ごと・感情ごとの数．本当がその感情のとき，その分類器をどれだけ頼ってよいか \\
水準／相性 & 信頼度の行の平均／平均からの感情ごとの差 \\
残り方 & ソースの外で測った信頼度 ÷ ソースの中で測った信頼度 \\
戻す強さ n0 & K 本で測った水準を，戻し先へどれだけ近づけるか．ラベルの本数で数える \\
切片 & 感情ごとに点数へ足す数．答えが特定の感情に片寄るのを直す \\
仮のターゲット & ソースの 1 つをターゲットに見立てたもの．全ラベルがあるので答え合わせができる \\
下調べ／評価の回 & 仮のターゲットで手法の形と設定を決める計算／決めた手法を本来のターゲットで 1 回だけ測る計算 \\
疑似ラベル & 人ではなく分類器が付けたラベル \\
UAR & 感情ごとの正解率の平均．でたらめなら 25 点 \\
基準 & 区間の下端が 0 より上で，差が 1 点以上．満たした条件に $\uparrow$ を付ける \\
\bottomrule
\end{tabular}
\end{center}

\section{結果ファイルの名前との対応}
\begin{center}
\small
\begin{tabular}{>{\raggedright}p{4.2cm}p{11.2cm}}
\toprule
名前 & 意味 \\
\midrule
\texttt{prod} & 積 \\
\texttt{M\_F6\_mixsel} & 提案（水準は K 本＋相性はソース．設定は別のコーパスで選ぶ） \\
\texttt{M\_F6\_K} & K 本だけで測った表 \\
\texttt{M\_F6\_ext} & ソースだけで測った表 \\
\texttt{M\_F6\_orc} & 全ラベルで測った表（参照） \\
\texttt{M\_F6n\_mixsel} & 提案から「どの感情でどのモダリティが強いか」の成分を消した版 \\
\texttt{none\_AVT}，\texttt{none\_AV} & K 本だけで学んだ融合の分類器（声・顔・言葉／声・顔） \\
\texttt{adj} & ラベルを使わずに切片をずらした答え（この資料の数はすべてこの答えの値） \\
\bottomrule
\end{tabular}
\end{center}
結果の置き場：4 コーパスの下調べは \texttt{results/dai16\_probe/}・\texttt{results/dai16\_probe\_k10/}，6 コーパスは \texttt{results/dai18\_probe/}，評価の回は \texttt{results/dai20\_eval/}．

\end{document}
